\section{Теорема Дедекинда о независимости характеров}

В этом и следующих трёх разделах один простой инструмент~--- линейная независимость различных гомоморфизмов~---
даёт три важных результата: теорему о нормальном базисе, теорему Гильберта~90 и описание циклических расширений.

\begin{definition}
Пусть $G$~--- группа, $F$~--- поле. \emph{Характер} группы $G$ со значениями в $F$~--- это гомоморфизм групп
$\chi\colon G\to F^\times$. Характеры~--- функции $G\to F$, поэтому их можно складывать и умножать на элементы $F$
поточечно: они лежат в $F$-векторном пространстве всех функций $G\to F$, и имеет смысл говорить об их линейной
зависимости над $F$.
\end{definition}

\begin{theorem}[5.14: Дедекинд]
Пусть $F$~--- поле, $G$~--- группа. Любой конечный набор $\{\chi_1,\dots,\chi_m\}$ \emph{различных} гомоморфизмов
$G\to F^\times$ линейно независим над $F$, т.\,е.
\[
a_1\chi_1+\dots+a_m\chi_m=0\ \ (\text{как функция } G\to F)\quad\Longrightarrow\quad a_1=\dots=a_m=0.
\]
\end{theorem}

Идея: если между характерами есть линейное соотношение, то, сдвигая аргумент, получаем второе соотношение; вычитая
из него первое, умноженное на подходящую константу, исключаем один из характеров и получаем более короткое соотношение.

\begin{proof}
Индукция по $m$. При $m=1$: если $a_1\chi_1=0$, то $a_1=a_1\chi_1(1)=0$, так как $\chi_1(1)=1$.

Пусть $m\ge2$ и утверждение верно для наборов из $m-1$ различных характеров. Пусть
\begin{equation}\label{eq:s4-ded1}
a_1\chi_1(x)+a_2\chi_2(x)+\dots+a_m\chi_m(x)=0\quad\text{для всех } x\in G.
\end{equation}
Так как $\chi_1\ne\chi_2$, найдётся $h\in G$ с $\chi_1(h)\ne\chi_2(h)$. Подставим в \eqref{eq:s4-ded1} $hx$ вместо $x$ и
воспользуемся тем, что $\chi_i(hx)=\chi_i(h)\chi_i(x)$:
\[
a_1\chi_1(h)\chi_1(x)+a_2\chi_2(h)\chi_2(x)+\dots+a_m\chi_m(h)\chi_m(x)=0\quad\text{для всех } x\in G.
\]
Вычтем отсюда \eqref{eq:s4-ded1}, умноженное на $\chi_1(h)$; первое слагаемое сократится:
\[
a_2'\chi_2+\dots+a_m'\chi_m=0,\qquad a_i'=a_i\bigl(\chi_i(h)-\chi_1(h)\bigr).
\]
По предположению индукции все $a_i'=0$. Так как $\chi_2(h)-\chi_1(h)\ne0$, получаем $a_2=0$. Теперь \eqref{eq:s4-ded1}
превращается в соотношение $a_1\chi_1+a_3\chi_3+\dots+a_m\chi_m=0$ между $m-1$ различными характерами, и по
предположению индукции остальные $a_i$ тоже равны нулю.
\end{proof}

Различность характеров существенна: при $\chi_1=\chi_2$ соотношение $\chi_1-\chi_2=0$ нетривиально. В явном виде она
используется при выборе $h$ (подсемейства, к которым применяется предположение индукции, состоят из различных
характеров автоматически).

\begin{remark}[5.14a]
Слово <<характер>> употребляется и в теории представлений: если $\rho\colon G\to GL_n(F)$~--- представление, то его
характером называют функцию $g\mapsto\operatorname{tr}\rho(g)$. При $n=1$ имеем $GL_1(F)=F^\times$ и
$\operatorname{tr}\rho=\rho$, так что характеры из определения выше~--- это характеры одномерных представлений. При $n>1$
функция $g\mapsto\operatorname{tr}\rho(g)$, вообще говоря, не гомоморфизм (в характеристике $0$~--- никогда:
$\operatorname{tr}\rho(1)=n\ne1$), и теорема \rf{5.14} о ней ничего не говорит.
\end{remark}

\begin{example}[5.14b]
(a) Пусть $G=\Z$, $a\in F^\times$ и $\chi_a(k)=a^k$. Это характер, и при $a\ne b$ характеры $\chi_a$, $\chi_b$
различны. По \rf{5.14} последовательности $(a_1^k)_{k\in\Z},\dots,(a_m^k)_{k\in\Z}$ с попарно различными $a_i$
линейно независимы над $F$ (это видно и из определителя Вандермонда: уже их отрезки при $k=0,\dots,m-1$ независимы).

(b) Пусть $G=(\R,+)$ и $\chi_\lambda(t)=e^{\lambda t}\in\R^\times$. Так как $e^{\lambda(s+t)}=e^{\lambda s}e^{\lambda t}$,
это характеры, попарно различные при различных $\lambda$. По \rf{5.14} экспоненты $e^{\lambda_1t},\dots,e^{\lambda_mt}$
с попарно различными $\lambda_i$ линейно независимы над $\R$~--- известный из теории дифференциальных уравнений факт.
\end{example}

\begin{corollary}[5.15]
Пусть $F$ и $E$~--- поля, и $\sigma_1,\dots,\sigma_m$~--- различные гомоморфизмы $F\to E$. Тогда
$\sigma_1,\dots,\sigma_m$ линейно независимы над $E$: если $c_1,\dots,c_m\in E$ и
$c_1\sigma_1(x)+\dots+c_m\sigma_m(x)=0$ для всех $x\in F$, то $c_1=\dots=c_m=0$.
\end{corollary}

\begin{proof}
Гомоморфизм полей инъективен, поэтому $\sigma_i(F^\times)\subset E^\times$, и $\chi_i=\sigma_i|_{F^\times}$~---
гомоморфизм групп $F^\times\to E^\times$. Эти характеры попарно различны: все $\sigma_i$ переводят $0$ в $0$, так что
если $\sigma_i$ и $\sigma_j$ совпадают на $F^\times$, то они совпадают всюду. Если $\sum c_i\sigma_i=0$ на $F$, то
$\sum c_i\chi_i=0$ на $F^\times$, и по \rf{5.14} (применённой к группе $F^\times$ и полю $E$) все $c_i=0$. (Значения в
нуле ничего не добавляют: $\sum c_i\sigma_i(0)=0$ при любых $c_i$.)
\end{proof}

\begin{remark}[5.15a]
Чаще всего \rf{5.15} применяют в такой форме. Пусть $E/F$~--- конечное расширение Галуа с группой $G$. Элементы $G$~---
различные гомоморфизмы $E\to E$, поэтому \emph{для любых $c_\sigma\in E$, не всех равных нулю, отображение
$\sum_{\sigma\in G}c_\sigma\sigma\colon E\to E$ не нулевое}, т.\,е.\ $\sum_\sigma c_\sigma\sigma(x)\ne0$ для некоторого
$x\in E$. Например, при всех $c_\sigma=1$ получаем: \emph{след} $\operatorname{Tr}_{E/F}=\sum_{\sigma\in G}\sigma\colon
E\to E$ не равен тождественно нулю. Так как $\tau\circ\operatorname{Tr}_{E/F}=\operatorname{Tr}_{E/F}$ для $\tau\in G$,
значения следа лежат в $E^G=F$, и $\operatorname{Tr}_{E/F}\colon E\to F$~--- ненулевое, а значит, сюръективное
$F$-линейное отображение. Так же $1,\sigma,\dots,\sigma^{n-1}$ линейно независимы над $E$, если $\sigma$ порождает
циклическую группу Галуа порядка $n$. Эти факты используются в \rf{5.18}, \rf{5.23} и \rf{5.27}.
\end{remark}

\begin{corollary}[5.16]
Пусть $E/F$~--- конечное сепарабельное расширение степени $m$, $\alpha_1,\dots,\alpha_m$~--- базис $E$ над $F$, и
$\sigma_1,\dots,\sigma_m$~--- различные $F$-гомоморфизмы $E$ в поле $\Om$. Тогда матрица $(\sigma_i\alpha_j)_{i,j}$ размера
$m\times m$ с элементами из $\Om$ обратима.
\end{corollary}

\begin{proof}
Пусть матрица вырождена. Тогда её строки линейно зависимы над $\Om$: найдутся $c_1,\dots,c_m\in\Om$, не все равные нулю,
с
\[
\sum_{i=1}^m c_i\,\sigma_i(\alpha_j)=0\quad\text{для всех } j=1,\dots,m.
\]
Отображение $\sum_ic_i\sigma_i\colon E\to\Om$ линейно над $F$ (каждое $\sigma_i$ тождественно на $F$), и оно обращается в
нуль на базисе $\alpha_1,\dots,\alpha_m$. Значит, $\sum_ic_i\sigma_i(\alpha)=0$ для всех $\alpha\in E$, что противоречит
\rf{5.15}.
\end{proof}

\begin{remark}[5.16a: где нужна сепарабельность]
В самом доказательстве сепарабельность не используется: рассуждение годится для любых $m$ различных
$F$-гомоморфизмов $E\to\Om$. Она нужна, чтобы условие было выполнимо, т.\,е.\ чтобы $m=[E:F]$ различных $F$-гомоморфизмов
вообще существовали. По \rf{2.14} их не больше $[E:F]$. Если $E/F$ сепарабельно, то $E$ содержится в некотором
расширении Галуа $K/F$ (\rf{3.12}), и если $\Om\supset K$, то ограничения на $E$ элементов $\Gal(K/F)$ дают
$(\Gal(K/F):\Gal(K/E))=[E:F]$ различных $F$-гомоморфизмов $E\to\Om$ (основная теорема \rf{3.17}). Для несепарабельного
расширения это не так: если $F=\F_p(t)$ и $E=F[t^{1/p}]$, то $X^p-t=(X-t^{1/p})^p$ в $E[X]$, поэтому любой
$F$-гомоморфизм $E\to\Om$ переводит $t^{1/p}$ в единственный корень $X^p-t$ в $\Om$, и такой гомоморфизм всего один, хотя
$[E:F]=p$. Тем же рассуждением для $k<m$ различных $F$-гомоморфизмов $\sigma_1,\dots,\sigma_k$ получаем, что матрица
$(\sigma_i\alpha_j)$ размера $k\times m$ имеет ранг $k$.
\end{remark}

\begin{remark}[5.16b]
Пусть $\sigma_1,\dots,\sigma_m$~--- все $F$-гомоморфизмы сепарабельного расширения $E$ в алгебраическое замыкание $\Om$
поля $F$, и $A=(\sigma_i\alpha_j)$. Тогда $(A^{\mathsf T}A)_{jk}=\sum_i\sigma_i(\alpha_j)\sigma_i(\alpha_k)
=\sum_i\sigma_i(\alpha_j\alpha_k)=\operatorname{Tr}_{E/F}(\alpha_j\alpha_k)$ по \rf{5.47a}(a). Значит,
$\det\bigl(\operatorname{Tr}_{E/F}(\alpha_j\alpha_k)\bigr)=(\det A)^2\ne0$: для сепарабельного расширения билинейная
форма $(x,y)\mapsto\operatorname{Tr}_{E/F}(xy)$ на $E$ невырождена. В таком виде \rf{5.16} используется в теории
чисел (дискриминант базиса).
\end{remark}
